% Standalone problem document, generated from the adjacent README.md.
% Compile with XeLaTeX. Mathematical content is maintained in Markdown.
\documentclass[11pt,a4paper]{article}
\usepackage[fontsize=10.5pt]{scrextend}
\usepackage[margin=22mm,headheight=15pt,headsep=9mm,footskip=11mm]{geometry}
\usepackage{fontspec}
\setmainfont{texgyrepagella}[Extension=.otf,UprightFont=*-regular,BoldFont=*-bold,ItalicFont=*-italic,BoldItalicFont=*-bolditalic]
\setsansfont{texgyreheros}[Extension=.otf,UprightFont=*-regular,BoldFont=*-bold,ItalicFont=*-italic,BoldItalicFont=*-bolditalic]
\setmonofont{texgyrecursor}[Extension=.otf,UprightFont=*-regular,BoldFont=*-bold,ItalicFont=*-italic,BoldItalicFont=*-bolditalic,Scale=MatchLowercase]
\usepackage{amsmath,amssymb}
\usepackage{unicode-math}
\setmathfont{texgyrepagella-math.otf}
\usepackage{xcolor,microtype,enumitem,needspace,fancyhdr}
\usepackage{bookmark}
\definecolor{ink}{HTML}{17344B}
\definecolor{accent}{HTML}{197D80}
\definecolor{muted}{HTML}{596773}
\hypersetup{colorlinks=true,linkcolor=accent,urlcolor=accent,pdftitle={MF-08 - NP-hardness
of unrestricted static output-feedback
stabilization},pdfauthor={Open Problems in Numerical Linear Algebra}}
\urlstyle{same}
\setlength{\parindent}{0pt}
\setlength{\parskip}{5pt plus 1pt minus 1pt}
\linespread{1.04}
\setlength{\emergencystretch}{3em}
\widowpenalty=10000
\clubpenalty=10000
\displaywidowpenalty=10000
\allowdisplaybreaks[1]
\setlist{nosep,leftmargin=1.5em}
\providecommand{\tightlist}{\setlength{\itemsep}{0pt}\setlength{\parskip}{0pt}}
\providecommand{\pandocbounded}[1]{#1}
\pagestyle{fancy}
\fancyhf{}
\fancyhead[L]{\sffamily\footnotesize\color{muted}OPEN PROBLEMS IN NUMERICAL LINEAR ALGEBRA}
\fancyhead[R]{\sffamily\bfseries\footnotesize\color{accent}MF-08}
\fancyfoot[L]{\parbox[b]{0.9\textwidth}{\sffamily\fontsize{8}{10}\selectfont\color{muted}Editorial ratings; a bounded search cannot certify that a problem remains unsolved.\\Literature check: 2026-09-10}}
\fancyfoot[R]{\sffamily\footnotesize\color{muted}\thepage}
\renewcommand{\headrulewidth}{0.3pt}
\renewcommand{\headrule}{\hbox to\headwidth{\color{accent}\leaders\hrule height \headrulewidth\hfill}}
\makeatletter
\renewcommand{\section}{\Needspace{5\baselineskip}\@startsection{section}{1}{0pt}{12pt plus 2pt minus 2pt}{4pt}{\sffamily\bfseries\large\color{ink}}}
\renewcommand{\subsection}{\Needspace{4\baselineskip}\@startsection{subsection}{2}{0pt}{9pt plus 1pt minus 1pt}{3pt}{\sffamily\bfseries\normalsize\color{ink}}}
\makeatother
\setcounter{secnumdepth}{0}
\begin{document}
{\sffamily\small\color{muted}Matrix functions and stability\par}
\vspace{3pt}
{\raggedright\hyphenpenalty=10000\exhyphenpenalty=10000\sffamily\bfseries\fontsize{21}{24}\selectfont\color{ink}NP-hardness
of unrestricted static output-feedback stabilization\par}
\vspace{7pt}
{\sffamily\small\textbf{Difficulty:} extreme\quad\textbf{Importance:} broadly
interesting\par}
{\sffamily\small\bfseries\color{accent}Status: Open\par}
\vspace{4pt}
{\color{accent}\hrule height 0.8pt}
\vspace{6pt}
\textbf{Rating rationale:} Extreme because this is a longstanding
complexity barrier for unrestricted feedback design; broad impact
connects control synthesis, optimization, and computational complexity.

\subsection{Problem statement}\label{problem-statement}

Consider the decision language whose inputs are rational matrices
\(A\in\mathbb Q^{n\times n}\), \(B\in\mathbb Q^{n\times m}\), and
\(C\in\mathbb Q^{p\times n}\), with positive dimensions included in the
input. The answer is yes precisely when there exists an unrestricted
real matrix \(K\in\mathbb R^{m\times p}\) for which

\[
\operatorname{Re}\lambda<0\qquad\text{for every }\lambda\in\sigma(A+BKC).
\]

Is this language NP-hard under polynomial-time many-one reductions in
the binary encoding of the rational input? The requested output is a
complexity theorem, not another numerical heuristic for finding \(K\).

\subsection{References and status
evidence}\label{references-and-status-evidence}

Minyue Fu,
\href{https://www.eng.newcastle.edu.au/~mf140/home/Papers/Fu_UTSC.pdf}{Two
Challenging Problems in Control Theory}, §3, pp.~3--7, defines strict
Hurwitz stabilization and distinguishes it from the known NP-hard
problem with entrywise bounds on \(K\) (§3.1). Gillis and Sharma,
\href{https://arxiv.org/pdf/2202.02618}{Solving matrix nearness problems
via Hamiltonian systems, matrix factorization, and optimization},
§3.5.2, pp.~50--51 (2022 manuscript), reiterates the unresolved
complexity of static output feedback. Its broader stability convention
should not replace the strict inequality above. Follow-up searches for
unrestricted stabilization NP-hardness found no resolving reduction;
claims about arbitrary bilinear matrix inequalities or prescribed pole
placement do not by themselves settle this language. No claim of NP
membership is made.

\subsection{Audit --- 2026-09-10}\label{audit-2026-09-10}

Rechecked
\href{https://www.eng.newcastle.edu.au/~mf140/home/Papers/Fu_UTSC.pdf}{Fu,
§3.1} and \href{https://arxiv.org/pdf/2202.02618}{Gillis--Sharma,
§3.5.2}. Searches for unrestricted static-output-feedback NP-hardness
found no applicable reduction. Bounded feedback gains, prescribed poles,
and general bilinear inequalities remain different decision problems;
later generic hardness descriptions do not establish this one.
\end{document}
